\section{Exact sequences and additive functors}
\label{sec:conventions}

\subsection{Coefficient modules and their Homs}

\begin{definition}\label{def:presentable-coefficients}
Let $\PrL$ have compactly generated presentable categories as objects
and compact-preserving left adjoints as morphisms. Its tensor product
represents functors preserving colimits separately
\citep[Section~4.8.1]{LurieHA}. We use the symmetric monoidal tower
\begin{equation}\label{eq:coefficient-tower}
 \nPr0=\Sp,\qquad
 \nPr{(r+1)}=\Mod_{\nPr r}(\PrL)\quad(r\geq0).
\end{equation}
Write $\one_r$ for the unit of $\nPr r$.
By \citep[Section~2.1]{Aoki}, the coefficient categories
$\nPr r$ are commutative algebra objects in $\PrL$.
In particular, $\nPr1$ consists of compactly generated
stable presentable categories.
\end{definition}

Ordinary Ind-completion gives the symmetric monoidal equivalence
\begin{equation}\label{eq:stable-ind-equivalence}
 \Ind:\Cat^{\mathrm{perf}}\xrightarrow{\sim}\nPr1,
\end{equation}
where $\Cat^{\mathrm{perf}}$ has small idempotent-complete stable
categories and exact functors. Its inverse takes compact objects
\citep{BGT,LurieHA}.

For $A\in\nPr{(r+1)}$, we write
$\PrHom_A(a,b)\in\nPr r$ for the
coefficient-enriched Hom, characterized by
\begin{equation}\label{eq:tensor-hom}
 \Map_{\nPr r}(v,\PrHom_A(a,b))
 \simeq\Map_A(v\odot a,b).
\end{equation}
Composition and units are supplied by the
coefficient enrichment \citep{GH,Heine}.
At $r=0$ this is the mapping spectrum.

We use the ordinary categories underlying the coefficient modules.
The coefficient-forgetting construction is
\cite[Remark~4.3]{ScholzeGeometry}, with the fixed
$\kappa=\omega$ formalism of
\cite[Definition~1.2 and Section~2.1]{Aoki}.

For $H\in\nPr r$, $r\geq1$,
\[
 \Map_{\nPr r}(\one_r,H)\simeq(H^\omega)^\simeq.
\]
Consequently, for $A\in\nPr{(r+1)}$ and $r\geq1$,
\[
 \Map_A(a,b)\simeq\bigl(\PrHom_A(a,b)^\omega\bigr)^\simeq.
\]
Thus an ordinary arrow in a module is a compact object of its
coefficient Hom. Subsequent local data are read in those coefficient
modules in the same way. We retain their full coefficient Homs.

\subsection{Exact sequences}

\begin{definition}\label{def:adjunction}
Let $H\in\nPr s$, $s\geq2$.
An internal adjunction $i\dashv r$ consists of
arrows $i:a\to b$ and $r:b\to a$, together
with a unit $\eta:1_a\to ri$, a counit
$\varepsilon:ir\to1_b$, and coherent triangle
data in the coefficient Homs.
Equivalently, these are the dual-pair data
of opposite off-diagonal matrices in
\[
 \mathcal B_H(a,b)=
 \prod_{u,v\in\{a,b\}}
 \PrHom_H(u,v)^\omega,
\]
with matrix multiplication given by composition
and finite biproducts.
\end{definition}


Let $H\in\nPr s$, $s\geq2$. Consider adjunctions
$p:b\rightleftarrows c:j$ and $q:b'\rightleftarrows c':k$,
together with arrows $F:b'\to b$ and $G:c'\to c$.
For a $2$-morphism $\beta:pF\Rightarrow Gq$, its right mate is
\[
 Fk\xrightarrow{\eta_p Fk}jpFk
   \xrightarrow{j\beta k}jGqk
   \xrightarrow{jG\varepsilon_q}jG,
\]
where $\eta_p:1_b\to jp$ and $\varepsilon_q:qk\to1_{c'}$
are the unit and counit of the respective adjunctions.
The Beck--Chevalley condition requires this mate to be invertible.



\begin{definition}\label{def:genuine}
For $H\in\nPr1$, an exact sequence is a specified cofiber square
\[
 \begin{tikzcd}
 a\ar[r]\ar[d]&b\ar[d]\\
 0\ar[r]&c.
 \end{tikzcd}
\]
For $H\in\nPr s$, $s\geq2$, an exact sequence
$a\xrightarrow{i}b\xrightarrow{p}c$ consists of $i\dashv r$,
$p\dashv j$, invertible unit $1_a\to ri$, invertible counit
$pj\to1_c$, $pi\simeq0$, and the recursively exact central sequence
\begin{equation}\label{eq:central-row}
 ir\xrightarrow{\varepsilon}1_b\xrightarrow{\eta}jp
       \quad\text{in }\PrHom_H(b,b).
\end{equation}
Its zero comparison is induced by the adjunctions and $pi\simeq0$.
All lower diagram data are retained. We abbreviate the invertible
unit and counit by $ri\simeq1_a$ and $pj\simeq1_c$.

\label{def:exact-sequence-category}
A morphism of complete exact-sequence diagrams has coherent
invertible comparisons for the forward arrows and their
specified right adjoints, compatible with the units and
counits. In particular, the corresponding
Beck--Chevalley mates are required to be invertible.

Write $\EqRow_s(H)$ for this concrete diagram object, including
strong maps of sequences and their coefficient Homs.
The notation $\EqRow_s^\omega(H)$ restricts its three vertices to
$H^\omega$.
\end{definition}

A cofiber square over its first arrow has a contractible space of
choices. With both successive arrows fixed, the specified square
is still part of the data. Adjunction gives $rj\simeq0$ as well.

\subsection{Diagrams lowering the coefficient index}


We write $\Fun^L_{\nPr r}(A,E)$ for the category of
coefficient-linear left adjoints, with strong transformations
and their coefficient Homs. The superscript $L,\omega$
additionally requires preservation of compact objects, while
$\Fun^\oplus$ denotes functors preserving zero objects and
finite biproducts.

Strong transformations have coherent invertible naturality
comparisons. Their coefficient Homs are given by enriched
ends and retain noninvertible intertwiners. The universal
properties below are understood at the level of coefficient
Homs. We do not assume that these functor categories are
compactly generated.



\begin{definition}\label{def:algebraic-diagrams}
For $A\in\nPr{(n+1)}$ and $E\in\nPr n$, define the category $\operatorname{Fun}_n^\omega(A,E)$ by recursion. At $n=1$ it is
$\Fun(A^\omega,E)$. At $n\geq2$ its objects assign $F(a)\in E$
to each $a\in A^\omega$ and local diagrams
\begin{equation}\label{eq:compact-algebraic-diagrams}
 F_{ab}\in\operatorname{Fun}_{n-1}^{\omega,c}
 \bigl(\PrHom_A(a,b),\PrHom_E(Fa,Fb)\bigr),
\end{equation}
with coherent invertible composition and unit comparisons.
Here $\operatorname{Fun}_m^{\omega,c}(B,D)$ is the full subobject
of $\operatorname{Fun}_m^\omega(B,D)$ whose outer values belong to
$D^\omega$. Thus $F(a)$ is arbitrary in $\operatorname{Fun}_n^\omega$,
whereas the local diagrams have compact local values at every index.
Transformation objects and their local Homs are defined by the
corresponding recursive ends.
The notation $\operatorname{Fun}_{n,\mathrm{filt}}(A,E)$ uses all
objects of $A$, with the same local condition, and requires its
ordinary object functor $A\to E$ to preserve filtered colimits.
\end{definition}


Informally, $\operatorname{Fun}_n^\omega(A,E)$ plays the role of
a higher functor category from the underlying $(\infty,n)$-category
of $A$ to that of $E$, with functorial data specified on
compact objects and subject to recursive compact-image
conditions on the coefficient Homs.




\begin{definition}\label{def:additive}
Let $\Add_n^\omega(A,E)$ be the full subobject of
$\operatorname{Fun}_n^\omega(A,E)$ on diagrams preserving zero and
exact sequences. Their images retain the specified
zero comparison and all adjunction data represented at the target
index. Define $\Add_{n,\mathrm{filt}}(A,E)$ analogously.
\end{definition}


We refer to $\Add_n^\omega(A,E)$ as the category of additive
invariants on compact objects, and to
$\Add_{n,\mathrm{filt}}(A,E)$ as the category of
filtered-colimit-preserving additive invariants.
Propositions~\ref{prop:filtered-extension}
and~\ref{prop:global-additive-extension}
show that restriction to compact objects
induces an equivalence between these categories.


Over a rigid symmetric monoidal stable base, these rows are the
split-exact sequences of \cite{BGT}. At the next coefficient level they are the
diagrams used by Mazel-Gee and Stern \citep{BGT,MGS}.
Section~\ref{sec:mapping} gives the normalizations.


